\section{Introduction}

\subsection{Motivation}
Enterprise documentation is overwhelmingly \emph{procedural}: it tells a user how to
accomplish a goal through an ordered sequence of steps. Adobe Experience Platform
(\aep) and Adobe Journey Optimizer (\ajo) ship thousands of documentation pages and
video tutorials describing how to build audiences, enrich profiles, design journeys, and
orchestrate campaigns. A retrieval or assistant system over this corpus must do more than
return topically similar passages---it must understand \emph{what comes before what}.
Concretely, three downstream capabilities motivate this survey:

\begin{enumerate}[leftmargin=1.4em,itemsep=0.15em]
  \item \textbf{Directional causal judgment.} Given two pieces of text
  (steps, phases, or doc fragments) $t_1$ and $t_2$, estimate
  $P(t_1 \text{ causally precedes / enables } t_2)$.
  \item \textbf{Follow-up question reranking.} Given a user question, order a slate of
  candidate follow-up questions by logical/causal progression, not just relevance.
  \item \textbf{Workflow ordering.} Given the unordered steps of a workflow, sort them
  into the correct ascending execution order.
\end{enumerate}

All three rest on a single primitive---a \emph{directional}, asymmetric notion of
precedence---that standard semantic embeddings, built around symmetric cosine
similarity, do not natively provide. We call a model that captures this primitive a
\emph{causal embedding model}, and we use the term deliberately broadly: it spans
asymmetric scoring heads over encoders, order/box/hyperbolic geometries, cross-encoder
reasoners, and small generative rerankers. The foundations section clarifies precisely
which rung of Pearl's causal ladder this task occupies (largely structural precedence,
not intervention).

\subsection{Why this is hard}
Three difficulties recur throughout the survey. First, \emph{symmetry}: the workhorse
similarity function $\cos(\mathbf{a},\mathbf{b})$ is symmetric, so a vanilla embedding
cannot, by construction, distinguish ``$A$ enables $B$'' from ``$B$ enables $A$.''
Second, \emph{causal direction is genuinely hard}: even frontier LLMs perform poorly on
pure causal-direction benchmarks, and supervised systems cap well below ceiling. Third,
and most consequentially for us, \emph{local correctness does not compose into global
correctness}: a model can be accurate on most individual pairs yet still produce a wholly
wrong global order, because the step that turns $O(n^2)$ noisy pairwise judgments into a
single ranking is fragile. In this project, a DeBERTa-v3 cross-encoder that orders a
6-step workflow nearly correctly and scores 21/30 on pairwise direction nonetheless
yields \textbf{0/30} correct \emph{full orderings} under a naive topological sort---a
gap this survey takes as its central diagnostic.

\subsection{Scope and contributions}
This document is both a \emph{survey} and a \emph{system-design study} grounded in a
real, ongoing project under a strict \textbf{sub-1B-parameter} budget and a
4$\times$A100 environment. Our contributions are:

\begin{itemize}[leftmargin=1.4em,itemsep=0.15em]
  \item A structured survey of the methods relevant to causal/directional embeddings,
  spanning eight areas: causal foundations, embedding/retrieval architectures, causal
  reasoning from text, asymmetric/order-structured geometry, reasoning rerankers and RL,
  sentence-ordering and rank aggregation, follow-up/conversational retrieval, and
  procedural-workflow understanding.
  \item An empirical characterization of the project's corpus, pairwise datasets, tier
  system, and follow-up data, with the baselines any new model must beat.
  \item A diagnosis that reframes the workflow-ordering failure as primarily an
  \emph{aggregation} and \emph{data} problem rather than a pairwise-model problem, and a
  concrete sub-1B system design that follows from it.
  \item A reference implementation (released alongside this survey) of robust ordering
  (minimum-feedback-arc-set, Bradley--Terry, Kemeny), score calibration, partial-order
  metrics, and a model-free follow-up reranker.
\end{itemize}

\subsection{Organization}
Section~2 establishes causal foundations and a working definition of ``causal
embedding.'' Section~3 surveys dense-retrieval architectures and modern embeddings.
Section~4 covers causal reasoning and direction from text. Section~5 develops asymmetric
and order-structured geometries---the most direct route to directionality. Section~6
surveys reasoning rerankers and RL for ranking under a small-model budget; Section~7
covers turning pairwise scores into a global order. Section~8 addresses follow-up and
conversational retrieval; Section~9, procedural-workflow understanding. Section~10
characterizes the data. Section~11 synthesizes a recommended system and roadmap, and
Section~12 concludes.
